\section{Limitations}
\label{sec:limitations}

Several constraints bound what can be concluded from this study.

\paragraph{Carbon is modelled, never metered.} Every energy and carbon figure in
this paper is computed from a fixed device power constant, a measured latency
and a grid-intensity constant (Section~\ref{sec:cost-accounting},
Table~\ref{tab:provenance}). No wall-power measurement was taken at any point.
The carbon comparisons between policies are therefore comparisons of modelled
quantities that are proportional to measured latency; they support statements
about which policy does more work, and they do not support any statement about
emissions of a deployed system. A metered replication on instrumented hardware
is the obvious next step and is not something this study attempted.


\paragraph{The carbon figures rest on two utilisation constants.} The energy
model reproduces, but only once the utilisation fractions are
supplied by hand. Across all \CarbonSequencesChecked{} antigens of the second
dataset, recomputing each per-antigen total from the recorded device
configuration, the measured latencies and the utilisation fractions of
Section~\ref{sec:cost-accounting} matches the stored values to within
\CarbonReconstructionError{}~gCO\textsubscript{2}e. The fractions are not in
the recorded configuration, so the stored totals cannot be rebuilt from that
configuration alone, and the released decomposition artifact says so. The
fractions themselves,
however, are literals in the benchmark code rather than entries in any
configuration file, and they dominate the result: the embodied term is almost
the whole of each per-request total.

That matters for how far the carbon conclusion travels. Of
\CarbonFractionSettingsTested{} alternative settings we evaluated, the ordering
survives in all but \CarbonFractionSettingsFlipping{}. A device pair with the
local tier at \CarbonFlipLocalFraction{} utilisation and the remote tier at
\CarbonFlipEdgeFraction{} would put the break-even routed fraction at
\CarbonFlipBreakEven{}, above the rate our policy actually routes at. Routing
would then be the greener option. Our finding therefore holds for the hardware
pair modelled here and for most reasonable variations on it. It remains a
statement about a cost model whose dominant parameters deserve to be measured
rather than assumed. A metered replication remains the right next
step.

\paragraph{The remote tier is co-located, not remote.} The latency figures
attribute a transport cost to routed antigens, but the benchmark executes both
tiers on one machine, and the two-node transport evidence in the artifact is an
HTTP probe, not a deployed remote model server. Real deployments would
add network variance, queueing and failure modes that
this harness does not model. For latency and for transmitted volume these
would, if anything, favour routing, not count against it, since
always-remote execution pays the network on every antigen and routing pays it
on roughly a quarter. What a real deployment would add against routing is of a
different kind: a dependency on the remote tier for exactly the antigens the
local model found hardest, and the failure and retry behaviour that comes with
it. We measure none of that.


\paragraph{The primary benchmark is small, and that bounds nearly everything.}
BP3C50ID's external split has \BepipredNumAntigens{} antigens. Resampling on
\BepipredNumAntigens{} units gives intervals wide enough that most quality
comparisons in this paper are not resolved, and every such comparison should be
read as "no effect detectable at this sample size" rather than as "no effect".
This applies to the routing results as much as to the fusion analysis: the
interval on the adaptive-split versus always-remote MCC difference spans roughly
$\pm 0.05$, which would not exclude a quality loss a practitioner would care
about. For the fusion analysis the IEDB subset at \NestedIedbSequences{}
antigens carries more weight on ROC-AUC and average precision, where it has no
tied folds. On the threshold-dependent metrics its informative-fold count is
closer to BP3C50ID's and we do not claim it is better powered there.

\paragraph{The language-model tier is the smallest available checkpoint.} The
PLM results use \texttt{esm2\_t6\_8M\_UR50D}, chosen for the hardware at hand.
Conclusions about the PLM family are conclusions about an 8M-parameter frozen
encoder and should not be read as characterising what language-model features
can contribute to this task.

\paragraph{One routing signal, one validation design.} Routing decisions come
from mean per-residue local uncertainty, selected on a single grouped validation
split per seed. Repeated grouped resampling of the validation split would give a
better picture of how stable the selected thresholds are; we vary the seed,
which moves the split, but we do not separate the two sources of variation.

\paragraph{The IEDB subset's sampling rule is not recorded.} Its source and
identity filter are established (Section~\ref{sec:datasets}), but how the
\IedbSubsetTest{} test and \IedbSubsetTrain{} training antigens were drawn from
the \IedbSourceSequences{} available is not captured in any script we could
find. The selection is therefore reproducible in its inputs but not in its
sampling. We state this rather than reconstruct a rule after the fact, since a
plausible-looking rule invented now would be indistinguishable in the write-up
from one that was actually followed, which is the failure this paper is about.
What we can do is pin down the set: the released manifest lists every test
antigen by accession, so the evaluation subset is identified exactly even
though the draw that produced it is not. The consequence for the reader is
that the IEDB arm should be treated as a second dataset of opportunity rather
than as a designed replication, and we do not rest the paper's conclusions on
it.
